% TOP_PROBLEM: {"id": 30002538, "title": "CAT(0) Actions of Artin Groups", "category_id": 17, "category": {"id": 17, "name": "group_theory", "display_name": "Group Theory", "description": "Problems about groups, group actions, representations, and related algebraic structures.", "slug": "group-theory", "order_index": 17, "created_at": "2026-07-31T15:26:25.671Z"}, "status": "open"}
% FIRST_POSTED: 2026-09-27
% !TeX program = lualatex
% ATTEMPT_STATUS: unresolved
% Model/process: GPT_6_Astra_Ultra; autonomous mathematical attempt,
% primary-source audit, exact finite diagnostics, and adversarial review.
% Prepared 27 September 2026; no general solution or novelty claimed.
\documentclass[11pt]{article}
\usepackage[margin=25mm]{geometry}
\usepackage{fontspec}
\setmainfont{Latin Modern Roman}
\usepackage{amsmath,amssymb,mathtools}
\usepackage{unicode-math}
\setmathfont{Latin Modern Math}
\usepackage{xurl,hyperref}
\hypersetup{colorlinks=true,urlcolor=blue}
\setcounter{secnumdepth}{0}
\setlength{\parindent}{0pt}
\setlength{\parskip}{5pt}
\setlength{\emergencystretch}{3em}
\begin{document}
\section*{\texorpdfstring{CAT(0) Actions of Artin Groups}{CAT(0) Actions of Artin Groups}}\label{30002538}
\subsection{Definitions and mathematical statement}
Let $S$ be a finite set and $(m_{st})$ a Coxeter matrix with $m_{ss}=1$ and $m_{st}\in\{2,3,\ldots,\infty\}$ for $s\ne t$. Its Artin group is
\[
 A=\langle S\mid
 \underbrace{sts\cdots}_{m_{st}\text{ letters}}=
 \underbrace{tst\cdots}_{m_{st}\text{ letters}}
 \text{ whenever }m_{st}<\infty\rangle.
\]
Finite matrix means finitely many generators, not that the associated Coxeter group is finite. A geodesic metric space is CAT(0) when distances within every geodesic triangle are at most the corresponding Euclidean comparison distances.

Does every such Artin group act properly and cocompactly by isometries on a CAT(0) space, in Haettel's Conjecture 1.1 convention? Properness requires only finitely many group translates of a compact set to meet that set, and cocompactness requires compact quotient. No particular standard complex, cube-complex structure, or faithful linear representation is prescribed as the construction.
\par\smallskip\textit{Formulation and concept references:} \href{https://www.numdam.org/item/10.5802/aif.3524.pdf}{[S1]}.

\subsection{Short English statement}
Does every finitely generated Artin group admit a proper cocompact isometric action on some CAT(0) space?
\subsection{Research attempt}

\paragraph{Scope, source audit, and strategy (27 September 2026).}
The quantifier is over every finite generating set and every permitted matrix;
the space and its dimension may depend on that matrix. The empty generating
set gives the trivial group acting on a point. Finite rank does not mean
spherical type. Throughout, trees have unit-length edges and products carry
the Euclidean product metric.

Haettel's Conjecture 1.1 in [S1] is exactly the geometric-action assertion
above. His Theorem 1.4 supplies compact locally CAT(0) models when every
finite label is at least five. This is a proved subcase, not an argument
covering the labels two, three, and four simultaneously. His Section 2 uses
dihedral groups as building blocks; the calculation below independently
spells out their tree-times-line action, without claiming novelty for that
known subcase. The accessible sources checked here do not establish the
all-matrix conjecture. A search for recent work was a scope check, not an
exhaustive certification that no further special cases exist.

The main attack is to start with explicit rank-two models, then ask what
compatibility would allow them to assemble. A second benchmark is the
right-angled cubical construction, whose local curvature can be checked
combinatorially. The counterexample track tests the simplest attempted
assembly: putting a Euclidean polygon on each defining relation. It fails
already for a group which does have a CAT(0) action. Consequently a failed
presentation complex cannot itself disprove the group conjecture.

\paragraph{An action lemma with all properness hypotheses exposed.}
\textbf{Lemma 500.1.}
Suppose a group $G$ acts without inversions on a locally finite tree $T$,
with finite quotient. Suppose there is a homomorphism $h:G\to\mathbb R$
whose restriction to every vertex stabilizer is injective with discrete
image. Finally suppose there is $z\in G$ fixing $T$ pointwise with
$h(z)\ne0$. Then
\[
 g\cdot(p,r)=(gp,r+h(g))
\]
is a proper cocompact isometric action on the CAT(0) space
$T\times\mathbb R$.

\emph{Proof.}
The action law follows from additivity of $h$. A metric tree is CAT(0),
and the squared distances in a product add; the CAT(0) comparison inequality
therefore passes to its Euclidean product with a line. The product is proper
as a metric space because the tree is locally finite.

Fix a vertex $v$. For any bound $R$, there are finitely many vertices $w$
with $d(v,w)\le R$. For each such $w$ in the orbit choose $g_wv=w$.
Every $g$ satisfying $gv=w$ is $g_wk$, with $k\in G_v$.
The further condition $|h(g)|\le R$ bounds $h(k)$ in an interval of finite
length. Discreteness and injectivity on $G_v$ leave finitely many $k$.
Thus there are finitely many group elements moving $(v,0)$ by at most any
fixed distance. If $gK\cap K\ne\varnothing$ for a compact $K$, choose a
ball of radius $R_K$ about $(v,0)$ containing $K$; then
$d((v,0),g(v,0))\le2R_K$. This proves the stated compact-set properness.

A finite union $Q$ of closed edges and vertices meets every $G$-orbit in
$T$. Move $p$ into $Q$ using an element $g$, and then use an integral power
of $z$ to move the second coordinate into $[0,|h(z)|]$. The latter operation
leaves the first coordinate unchanged. Hence the compact set
$Q\times[0,|h(z)|]$ meets every orbit. This proves cocompactness.
\hfill$\square$

The discrete-image condition cannot be replaced merely by a nonzero
homomorphism somewhere in $G$. Infinite vertex stabilizers must be controlled
where they occur. Likewise, the action on $T$ alone is not proper in the
examples that follow; adjoining the line removes that exact defect.

\paragraph{Odd dihedral labels: a full algebra-to-geometry calculation.}
Let $m=2k+1\ge3$, put $u=st$, and write
\[
 G_m=\langle s,t\mid (st)^ks=(ts)^kt\rangle.
\]
Since $ts=s^{-1}us$ and $t=s^{-1}u$, the defining relation becomes
$u^ks=s^{-1}u^{k+1}$, equivalently $su^ks=u^{k+1}$.
Putting $v=u^ks$ gives
\[
 v^2=u^k(su^ks)=u^{2k+1}=u^m.
\]
Conversely, from $\langle u,v\mid v^2=u^m\rangle$ define
$s=u^{-k}v$ and $t=v^{-1}u^{k+1}$. Then $st=u$ and
$su^ks=u^{-k}v^2=u^{k+1}$, recovering the original relation.
These substitutions are inverse, so they prove the presentation isomorphism
rather than only a quotient map:
\begin{equation}\label{ra500:odd}
 G_m\cong\langle u,v\mid u^m=v^2\rangle
       =\langle u\rangle*_{\langle u^m\rangle=\langle v^2\rangle}
        \langle v\rangle.
\end{equation}

The Bass--Serre tree of this amalgam has vertex stabilizers conjugate to
$\langle u\rangle$ and $\langle v\rangle$, and valences $m$ and $2$.
It has one edge in the quotient. These facts follow directly from its coset
description: vertices are $G_m/\langle u\rangle$ and
$G_m/\langle v\rangle$, edges are $G_m/\langle u^m\rangle$, and an edge
$g\langle u^m\rangle$ joins the two vertices represented by $g$.
The normal form for an amalgam says a reduced alternating product of
non-edge-group factors is nontrivial; it rules out a nonbacktracking cycle
in this connected coset graph, so the graph is a tree.

The element $z=u^m=v^2$ commutes with both generators and belongs to the
edge group, hence fixes every vertex and edge. The homomorphism
\[
 h(u)=2,\qquad h(v)=m
\]
is well-defined because both sides of the relation have height $2m$.
Its restrictions to the vertex stabilizers are injective discrete maps,
also after conjugation, since $h(gag^{-1})=h(a)$. Moreover $h(z)=2m$.
Lemma 500.1 now proves a proper cocompact action. Notice that the recovered
original generators have $h(s)=h(t)=1$, as required for consistency with
the homogeneous Artin relation.

\paragraph{Even dihedral labels, including the commuting boundary case.}
For $m=2k\ge2$, again put $u=st$. The equality
$(st)^k=(ts)^k$ is precisely $u^k=s^{-1}u^ks$, so
\begin{equation}\label{ra500:even}
 G_{2k}\cong\langle u,s\mid su^ks^{-1}=u^k\rangle.
\end{equation}
The inverse substitution is $t=s^{-1}u$. This is an HNN extension of
$\langle u\rangle\cong\mathbb Z$ with both associated subgroups
$\langle u^k\rangle$, identified by the identity. Its coset tree has
vertex stabilizers conjugate to $\langle u\rangle$, valence $2k$, and a
single loop in the quotient. The tree property is the HNN normal-form
statement: a reduced word containing a stable letter cannot represent the
identity unless a cancellable associated-subgroup segment occurs.

Here $z=u^k$ is central and fixes the entire tree. Set $h(u)=2$ and
$h(s)=1$. Again the relation preserves height, $h$ is injective discrete
on each vertex stabilizer, and $h(z)=2k$. Lemma 500.1 applies.
When $k=1$ the tree is a line and this construction produces the expected
Euclidean-plane action of $G_2=\mathbb Z^2$; it is not necessary to exclude
this degenerate valence-two case. For $m=\infty$ the group is free on
$s,t$ and its ordinary Cayley tree already supplies a proper cocompact
action. A one-generator Artin group acts on $\mathbb R$ by integral
translations. Thus all groups of rank at most two have been covered.

The underlying Bass--Serre normal forms are the only group-action input
in these constructions; the corresponding dihedral presentations and
geometry are also recorded in [S1, Section 2]. They are not being used to
assert that arbitrary unions of these models are nonpositively curved.

\paragraph{A separate all-rank benchmark: right-angled matrices.}
Assume every off-diagonal entry is $2$ or $\infty$. Let $\Gamma$ be the
graph of the entries equal to $2$. Construct a finite cubical complex $X$
with one vertex, one oriented circle for each generator, and a coordinate
$q$-torus for each $q$-clique of $\Gamma$, glued on the corresponding
coordinate subtori. Its two-cells impose exactly the commuting relations;
cells of dimension at least three do not change the fundamental group.
Consequently $\pi_1X$ is the desired Artin group.

Here is the full combinatorial curvature check. A link vertex is a signed
generator $s^+$ or $s^-$. A collection spans a simplex precisely when its
underlying generators are distinct and form a clique of $\Gamma$.
Opposite signs of the same generator are not adjacent. Thus a pairwise
adjacent collection has distinct underlying generators which commute
pairwise, so the required torus exists and fills its link simplex. Every
clique in the link is therefore filled: the link is flag.

The cubical link criterion and the CAT(0) globalization theorem imply that
the universal cover $\widetilde X$ is CAT(0); the exact statement needed is
that a finite Euclidean cube complex with flag links has a CAT(0) universal
cover in the lifted length metric [E1, Section 2]. Its finite cell structure
ensures local finiteness and completeness of that lifted metric. Deck
transformations act properly, and the quotient $X$ is compact. This
rederives the known right-angled subcase with no restriction on rank.

The importance of the higher cubes is visible for three mutually commuting
generators. If only the three commutator squares are retained, the link has
an unfilled triangle whose edge lengths are $\pi/2$; the loop has length
$3\pi/2<2\pi$. Attaching the three-torus provides the missing simplex.
Having each rank-two square locally suitable is therefore insufficient
even in this particularly simple setting.

\paragraph{A concrete failed construction for a noncommuting relation.}
Take the presentation $\langle s,t\mid sts=tst\rangle$. Its standard
two-complex has a hexagonal cell with attaching word
\[
 stst^{-1}s^{-1}t^{-1}.
\]
If that cell is a regular Euclidean hexagon, every corner contributes a
link edge of length $2\pi/3$. A consecutive pair of letters $a,b$ contributes
an edge joining the outgoing directions $a^{-1}$ and $b$.
The corner $(s,t)$ and the distinct corner $(t^{-1},s^{-1})$ both join
$s^-$ to $t^+$. They are different link edges with the same endpoints.
Their union is a locally geodesic two-edge loop of length $4\pi/3$.
The link fails the CAT(1) condition, and this polygonal metric is not locally
CAT(0). This is an exact obstruction, not an impression based on a picture.

The same group is $G_3$, for which the tree-times-line construction above
works. Thus replacing the group-existence question by the assertion that
its most obvious presentation metric is CAT(0) would change the problem
and yield a false stronger assertion. Subdividing the hexagon without
changing the metric does not remove this short link loop.

\paragraph{Where assembly stops and what would suffice next.}
The Artin presentation is assembled from two-generator relations, but a
geometric amalgamation requires more than matching those relations.
One would need compact local models with compatible lengths for each shared
generator loop and a proof that every link in the assembled space is CAT(1),
including loops that traverse three or more different rank-two blocks.
Convex-gluing theorems apply to suitably embedded convex subsets in CAT(0)
spaces; an arbitrary immersed loop in a quotient is not such a subset.
The preceding two examples show specific ways in which this distinction
matters. Neither abstract presentation gluing nor the existence of a
rank-two action verifies these hypotheses.

A precise sufficient next lemma for this route would be: for every finite
Artin matrix one can choose finitely many compact piecewise Euclidean
blocks, attach them along the generator data to obtain the stated Artin
fundamental group, and verify the CAT(1) inequality for all resulting
links. The unsupplied part is the last compatibility assertion, particularly
when small and large labels interact. It is not proved by the rank-two
calculations. The XXL construction in [S1] gives one carefully controlled
successful assembly, which explains why its angle hypotheses cannot simply
be dropped.

\paragraph{Verification and outcome.}
The algebraic substitutions above were checked in both directions; the
height of the central translation is nonzero in each finite-label case.
Properness was checked on compact sets, not inferred solely from a trivial
point stabilizer. The $m=2$, $m=3$, $m=\infty$, rank-one, and empty-rank
cases were separated. A small exact combinatorial script in the local
research record enumerates the six corners of the braid hexagon and
confirms the duplicated link edges; it also checks the signed-link flag
description on all 1,100 graphs with at most five vertices. These finite tests
are diagnostics of the constructions, not evidence for all matrices.

\textbf{Outcome: UNRESOLVED, with complete reconstructions of known subcases.}
The present work proves the explicit rank-two geometric actions and the
right-angled construction, and refutes two tempting shortcuts for gluing.
It supplies neither an all-matrix construction nor a non-CAT(0) Artin group.
The first unresolved step is compatible local curvature control when the
rank-two pieces are assembled beyond these verified families. No novelty
or resolution of the catalogue conjecture is claimed.

\subsection{Sources}
\begin{itemize}
\item[S1]T. Haettel, XXL type Artin groups are CAT(0) and acylindrically hyperbolic, AIF 72 (2022), 2541–2555.
\url{https://www.numdam.org/item/10.5802/aif.3524.pdf}.
\textit{Formulation source.}
\item[E1]Ruth Charney, \emph{An introduction to right-angled Artin groups},
Geometriae Dedicata 125 (2007), 141--158, Section 2.4 (link criterion),
Section 2.6 and Theorem 2.6 (Salvetti complex).
\url{https://arxiv.org/pdf/math/0610668}.
\textit{Author manuscript inspected 27 September 2026.}
\end{itemize}
\end{document}
